\documentclass{beamer}
\usetheme{Darmstadt}
\useoutertheme[subsection=false, footline=authortitle]{miniframes}
\usecolortheme[RGB={44, 131, 82}]{structure}

\title{Research and training report for the academic year 2016/2017}
\subtitle{Merged mathematical models of the polluted atmosphere}
\author{Nora Juhasz}
\institute{Department of Information Engineering, Computer Science and Mathematics \\ University of L'Aquila \\ Supervisor: Prof. Donatella Donatelli}

\date{27th October, 2017}

\pgfdeclareimage[width = 10mm]{disimlogo}{LogoDisim.png}
\logo{\pgfuseimage{disimlogo}}

\begin{document}

\begin{frame}
\titlepage
\end{frame}

\section{Presentation of fellow's background}

\begin{frame}{Presentation of fellow's background}
  \begin{itemize}
  \item
    Master degree in 2013, University of Szeged, Hungary
  \item
    Starting date of Marie Curie fellowship at the University of L'Aquila: 1st of October, 2016
  \item
    The first training year at the University of L'Aquila has consisted of the following type of activities:
    \begin{itemize}
      \item Training
      \item Research
    \end{itemize}
  \end{itemize}
\end{frame}

\section{Training activities in the academic year 2016/2017}

\begin{frame}{Training activities in the academic year 2016/2017}

  \begin{itemize}
  \item
    \textbf{University of L'Aquila}: Numerical Analysis 2, High Performance Parallel Computing
  \item
    \textbf{Gran Sasso Science Institute}: Introduction to Finite Elements Method, Decay property for hyperbolic systems with dissipation, Advanced Topics on Numerical Methods, Introduction to the incompressible Naver-Stokes equation
  \item
    \textbf{Seminars, workshops}: PDEs Day (GSSI), School on Uncertainty Quantification for Hyperbolic Equations and Related Topics, CIME School 2017  - Mathematical Analysis of the Navier-Stokes Equations (Cetraro)
  \end{itemize}
\end{frame}

\section{Current state of the research project}

\begin{frame}{Modeling of the atmosphere}

  The beginning step of this research project was to collect sets of equations that describe well the motion of the atmosphere, firstly without considering the effect of pollution. The nature / source of these equations are the following:
  \begin{itemize}
  \item
    conservation of momentum, Navier-Stokes equation
  \item
    conservation of mass (continuity equation)
  \item
    thermodynamics equation (conservation of energy)
  \item
    equation of state
  \item
    equation describing the quantity of water vapor (ocean analogy: salinity)
  \end{itemize}

\end{frame}

\begin{frame}{Different approaches for the equations}

  Using these equations we can create three main different versions of sets of equations that describe the atmosphere:
  \begin{itemize}
  \item
    Using the classical $(x, y, z)$ cartesian coordinates,
  \item
    Introducing a terrain-influenced coordinate $\eta$ instead of $z$,
  \item
    Using the well-behaving, monotonous relationship between pressure and height and use $p$ as a vertical coordinate.
  
  \end{itemize}

\end{frame}

\begin{frame}{An example set of equations}

As an example, the equations of the atmosphere in $(x, y, p)$ coordinates:

\[ \frac{\partial \bold{v}}{\partial t} + \nabla_{\bold{v}}\bold{v} + \omega \frac{\partial \bold{v}}{\partial p} + 2 \Omega \sin \theta \bold{k} \times \bold{v} + \nabla \Phi -L_{\bold{v}} \bold{v} = F_\bold{v}  \]

\[ \frac{\partial \Phi}{\partial p} + \frac{RT}{p} = 0, \nabla \cdot \bold{v} + \frac{\partial \omega}{\partial p} = 0 \]

\[ \frac{\partial T}{\partial t} + \nabla_{\bold{v}} T + \omega \frac{\partial T}{\partial p} - \frac{R \bar{T}}{c_p p} \omega - L_T T = F_T\]

\[ \frac{\partial q}{\partial t} + \nabla_{\bold{v}}q + \omega \frac{\partial q}{\partial p} -L_q q = F_q \]

\[ p = R \rho T\]

\end{frame}

\begin{frame}{Equation describing pollution}

The species continuity equation (equations of chemical kinetics) has the following form:

\[ \frac{\partial C_j}{\partial t} + \nabla \cdot J_j = r_j,\]

where

\begin{itemize}
\item $C_j$ is the local molar concentration of species $j$, (time variations in the species concentration),
\item the components of $J_j$ represent the local molar flux, (local motion of the species), and
\item $r_j$ is a volume based formation rate.
\end{itemize}

\end{frame}

\begin{frame}{Syncronising}
To merge the equations of the atmosphere and the equation of pollution we use the following approach:
\begin{itemize}
\item we write the fluid equations of the mixture considered as a single fluid with denisty $\rho_0$, pressure $p$, temperature $T$, velocity $\bold{v}$,
\item we consider the equations describing the evolution of each mass fraction (or species concentration) $Y_i$ or ($C_i$). 
\item we add "connecting" / "glue" equations regarding the quantities of the species to sum up for the quantity presented in the atmosphere equation.
\end{itemize}
\end{frame}

\begin{frame}{Syncronising}
One of the issues about syncronising was the fact that the equations of the atmosphere and pollution are often not in the same coordinate system.
For solving this we could either
\begin{itemize}
\item use Jacobians in the equations, or
\item use the cartesian coordinate system for both systems
\end{itemize}
We have chosen the later, since because of the quasi-horizontal nature of the atmosphere we will perform dimension reduction. 
\end{frame}

\begin{frame}{Current state part 1}

\[ \frac{\partial \bold{v} }{\partial t} + \nabla_{\bold{v}}\bold{v} + w  \frac{\partial \bold{v} }{\partial z} + \frac{1}{\rho_0}\nabla p + 2 \Omega \sin \theta \bold{k} \times \bold{v}  -\mu_{\bold{v}}\Delta \bold{v}  -\nu_{\bold{v}} \frac{\partial ^ 2 \bold{v} }{\partial z^2} = 0  \]

\[ \frac{\partial p }{\partial z} = - \rho_0 g \]

\[ \frac{\partial \rho_0}{\partial t} + \rho_0 (\nabla \bold{v} + \frac{\partial w }{\partial z}) + \bold{v} \nabla \rho_0 + w \frac{\partial \rho_0}{\partial z} = 0 \]

\[ \frac{\partial T }{\partial t} + \nabla_{\bold{v}} T  + w  \frac{\partial T }{\partial z} - \mu_T \Delta T  -\nu_T \frac{\partial^2 T }{\partial z^2} -\frac{RT}{p} \frac{\mathrm{d}p}{\mathrm{d}t} = Q_T \]

\[ \frac{\partial q }{\partial t} + \nabla_{\bold{v}}q + w \frac{\partial q}{\partial z} - \mu_q \Delta q  - \nu_q \frac{\partial ^2 q }{\partial z^2} = 0 \]

\[ p = R \rho_0 T \]

\end{frame}

\begin{frame}{Current state part 2}

 \[ \frac{\partial \bar{\rho_i}}{\partial t}  + \nabla_{h} (\bar{\rho_i} \bar{V_h}) + \frac{\partial}{\partial z} (\bar{\rho_i} \bar{V_v}) + \frac{\partial}{\partial x} (-\bar{\rho}K_{11} \frac{\partial \bar{Y_i}}{\partial x}) + \]
 \[+ \frac{\partial}{\partial y} (-\bar{\rho}K_{22} \frac{\partial \bar{Y_i}}{\partial y}) + \frac{\partial}{\partial z} (- \bar{\rho} K_{33} \frac{\partial \bar{Y_i}}{\partial z}) = \]
 \[ = R_{\rho_i} (\bar{\rho_1}, \dots , \bar{\rho_N}) + Q_{\rho_i} + \frac{\partial (\bar{\rho_i})}{\partial t} \rvert_{cld} + \frac{\partial (\bar{\rho_i})}{\partial t} \rvert_{aero} + \frac{\partial (\bar{\rho_i})}{\partial t} \rvert_{ping},\]
 
 (which is a more refined version of the previous pollution equation)

\[ \rho_0 = \sum_{i =1}^{N} \rho_i, p = \sum_{i =1}^{N} p_i, (\bold{v}, w) = \sum_{i =1}^{N}\frac{\rho_i}{\rho_0} (\bold{v}, w)_i \]

\end{frame}

\begin{frame}{Outlook}
  \begin{itemize}
  \item Dimension reduction for the merged system. The 3D model is too complex to work with, on the other hand the phenomena can be seen as quasi-horizontal, which is why we perform dimension reduction for the merged system, the goal is to have a synchronized 2D system describing pollution and atmosphere.
  \item Solutions for the for the new, two-dimensional merged system.
  \item Numerical simulations of solutions.
  \end{itemize}
\end{frame}

\end{document}